\documentclass[../../../include/open-logic-section]{subfiles}

\begin{document}

\olfileid[hi]{sth}{replacement}{absinf}
\olsection{प्रतिस्थापन और ``निरपेक्ष अनंतता''}

अब हम \limofsize{} को पीछे छोड़कर प्रतिस्थापन का औचित्य देने वाले (आंतरिक)
प्रयासों के एक भिन्न परिवार का अन्वेषण करेंगे, जो समुच्चयों के चरणों में बनने
के विचार को गंभीरता से लेते हैं।

पुनरावर्ती प्रक्रिया की पहली रूपरेखा देते समय हमने कुछ सिद्धांत दिए जो हर
चरण पर होने वाली बात समझाते थे। ये \stageshier, \stagesord और \stagesacc
थे। बाद में हमने कुछ सिद्धांत जोड़े जो चरणों की संख्या के विषय में बताते थे:
\stagessucc{} ने बताया कि समुच्चय-निर्माण की प्रक्रिया कभी समाप्त नहीं होती,
और \stagesinf{} ने बताया कि प्रक्रिया किसी अनंतवें चरण से होकर जाती है।

यह सुझाव उचित है कि ये अंतिम दो सिद्धांत किसी अधिक व्यापक सिद्धांत से
निष्पन्न होते हैं, जैसे:
\begin{enumerate}
	\item[] \stagesinex. निरपेक्ष रूप से अनंततः अनेक चरण हैं; पदानुक्रम
	उतना ऊँचा है जितना संभवतः हो सकता है।
\end{enumerate}
स्पष्टतः यह अनौपचारिक सिद्धांत है। किंतु यदि यह समुच्चय की संचयी-पुनरावर्ती
अवधारणा से तुरंत \emph{निष्पन्न} न भी हो, तो निश्चित ही उसके \emph{अनुकूल}
लगता है। कम-से-कम \limofsize{} के विपरीत यह इस विचार को बनाए रखता है कि
समुच्चय चरण-दर-चरण बनते हैं।

अब आशा है कि \stagesinex{} का उपयोग करके प्रतिस्थापन का औचित्य दिया जाए।
देखें कि यह कैसे हो सकता है।

\olref[ordinals][idea]{sec} में हमने देखा कि ऐसी सुव्यवस्था बनाना सरल है जो
(सारतः) $\omega+\omega$ के समाकृतिक होनी चाहिए। दूसरे शब्दों में हम आसानी
से ऐसी चरण-दर-चरण पुनरावर्ती प्रक्रिया की कल्पना कर सकते हैं जिसका
क्रम-प्रकार (सारतः) $\omega+\omega$ है। अतः यदि हमने \stagesinex{} स्वीकार
किया है, तो निश्चित ही स्वीकार करना चाहिए कि पदानुक्रम का कम-से-कम
$\omega+\omega$वाँ चरण, अर्थात $V_{\omega+\omega}$, है; क्योंकि पदानुक्रम
निश्चित ही इतनी दूर तक चल \emph{सकता} है।

यह विचार इस प्रकार सामान्यीकृत होता है: किसी भी सुव्यवस्था के लिए
पुनरावर्ती पदानुक्रम बनाने की प्रक्रिया को कम-से-कम उस सुव्यवस्था जितनी दूर
चलना चाहिए। और \olref[ordinals][ordtype]{thmOrdinalRepresentation} को
\emph{अभिगृहीत} मानकर इसकी गारंटी दी जा सकती है। यह बताएगा कि प्रत्येक
सुव्यवस्था किसी वॉन नोयमान क्रमसूचक के समाकृतिक है। चूँकि प्रत्येक वॉन
नोयमान क्रमसूचक अपनी रैंक के बराबर होगा,
\olref[ordinals][ordtype]{thmOrdinalRepresentation} हमें बताएगा कि जब भी हम
अपने समुच्चय सिद्धांत में कोई सुव्यवस्था वर्णित कर सकें, समुच्चय-निर्माण की
पुनरावर्ती प्रक्रिया को उस सुव्यवस्था से आगे निकलना होगा।

यह विचार निश्चित ही \stagesinex{} का उपप्रमेय जैसा लगता है। दुर्भाग्य से
यदि हमारा उद्देश्य इस विचार से प्रतिस्थापन निकालना है, तो एक सरल तकनीकी बाधा
सामने आती है: प्रतिस्थापन
\olref[ordinals][ordtype]{thmOrdinalRepresentation} से कड़ाई से अधिक प्रबल
है। (यह प्रेक्षण \citet[\S13.2]{Potter2004} का है; हम इसे
\olref[sth][replacement][finiteaxiomatizability]{sec} में सिद्ध करेंगे।)

निष्कर्ष है कि यदि हम \stagesinex{} को इस प्रकार समझना चाहते हैं कि उससे
प्रतिस्थापन मिले, तो वह \emph{केवल} यह नहीं कह सकता कि पदानुक्रम किसी भी
सुव्यवस्था से आगे निकलता है। उसे इससे अधिक प्रबल दावा करना होगा। इस उद्देश्य
से \cite{Shoenfield:AST} ने विचार का यह अत्यंत स्वाभाविक प्रबलीकरण सुझाया:
पदानुक्रम किसी भी समुच्चय के साथ \emph{सह-अंतिम} नहीं है।\footnote{प्रतीत
होता है कि ग्योडेल ने समान विचार सुझाया था; \citet[p.~223]{Potter2004}
देखें। ग्योडेल और \citeauthor{Shoenfield:AST} की चर्चा के लिए
\citet[90--5]{Incurvati2020} देखें।} कुछ अधिक विस्तार से: यदि $\tau$ ऐसा
प्रतिचित्रण है जो समुच्चयों को पदानुक्रम के चरणों पर भेजता है, तो $\tau$ के
अधीन किसी भी समुच्चय $A$ का प्रतिबिंब पदानुक्रम को समाप्त नहीं करता। दूसरे
शब्दों में (योजनाबद्ध रूप में):
\begin{enumerate}
	\item[] \stagescofin. यदि $A$ समुच्चय है और प्रत्येक $x \in A$ के लिए
	$\tau(x)$ कोई चरण है, तो ऐसा चरण है जो प्रत्येक $x \in A$ के लिए
	$\tau(x)$ के बाद आता है।
\end{enumerate}
स्पष्ट है कि $\ZF$, \stagescofin{} के उपयुक्त औपचारिक रूप को सिद्ध करता
है। इसके विपरीत हम अनौपचारिक तर्क दे सकते हैं कि \stagescofin{} प्रतिस्थापन
का औचित्य देता है।\footnote{\stagescofin{} के किसी औपचारिक रूप से
प्रतिस्थापन सिद्ध करना अधिक कठिन होगा, क्योंकि अकेला $\Z$ चरणों को परिभाषित
करने के लिए पर्याप्त प्रबल नहीं; अतः यह स्पष्ट नहीं कि \stagescofin{} को
औपचारिक कैसे बनाया जाए। एक विकल्प $\LT$ के किसी विस्तार में काम करना है,
जैसा \olref[sth][replacement][extrinsic]{sec} में चर्चा की गई।} मान लें
$(\forall x \in A)\lexists![y][\phi(x,y)]$। तब प्रत्येक $x \in A$ के लिए
$\sigma(x)$ को ऐसा $y$ मानें कि $\phi(x,y)$, और $\tau(x)$ को वह चरण मानें
जिस पर $\sigma(x)$ पहली बार बना। \stagescofin{} से ऐसा चरण $V$ है कि
$(\forall x \in A)\tau(x)\in V$। अब चूँकि प्रत्येक $\tau(x) \in V$ और
$\sigma(x) \subseteq \tau(x)$, इसलिए पृथक्करण से
$\Setabs{y \in V}{(\exists x \in A)\sigma(x) = y}
= \Setabs{y}{(\exists x \in A)\phi(x,y)}$ प्राप्त कर सकते हैं।

\begin{prob}
	$\ZF$ के भीतर \stagescofin{} को औपचारिक रूप दीजिए।
\end{prob}

अतः \stagescofin{} प्रतिस्थापन का औचित्य सिद्ध करता है। और कम-से-कम यह
विश्वसनीय है कि \stagesinex{}, \stagescofin{} का औचित्य सिद्ध करता है। मान
लें \stagescofin{} विफल होता है। तब पदानुक्रम किसी समुच्चय~$A$ के साथ
सह-अंतिम है; अर्थात हमारे पास ऐसा प्रतिचित्रण $\tau$ है कि किसी भी चरण~$S$
के लिए कोई $x \in A$ ऐसा है कि $S \in \tau(x)$। उस स्थिति में पदानुक्रम की
कथित ``निरपेक्ष अनंतता'' को पकड़ने का साधन हमारे पास है: $A$ पर लागू
$\tau$ का परास उसे \emph{समाप्त} कर देता है। और यह इस विचार से समझौता करता
है कि पदानुक्रम ``निरपेक्ष रूप से अनंत'' है। प्रतिधनात्मक रूप में:
\stagesinex{} से \stagescofin{} निष्पन्न होता है, जो बदले में प्रतिस्थापन का
औचित्य देता है।

यह प्रतिस्थापन का आंतरिक औचित्य देने का सचमुच आशाजनक प्रयास है। किंतु अंततः
यह सफल होता है या नहीं, इसका निर्णय हमें आप पर छोड़ना होगा।

\end{document}
